\chapter{EBM and the alternative model families}
\label{ch:ebm}

\section{The Explainable Boosting Machine}

EBM is a generalized additive model with bounded interactions, fitted by
cyclic gradient boosting on one feature at a time:
\begin{equation}
g\!\left(\mathbb{E}[y \mid \mathbf{x}]\right)
= \beta_0
+ \sum_{j=1}^{p} f_j(x_j)
+ \sum_{(j,k) \in \mathcal{I}} f_{jk}(x_j, x_k),
\qquad |\mathcal{I}| = 10 .
\end{equation}
Because each $f_j$ is a one-dimensional step function of a single feature,
the fitted model is directly readable: one can plot the learned effect of
storage temperature on shelf life and inspect it. This interpretability is
why it is in the production candidate set.

Its production configuration is \texttt{interactions=10},
\texttt{max\_bins=256}, \texttt{learning\_rate=0.02},
\texttt{max\_rounds=3000}, \texttt{min\_samples\_leaf=4},
\texttt{outer\_bags=8}.

\subsection{EBM's original results}

From \Cref{tab:tab_production_v7}, EBM is consistently the weakest of the
four production families under the original protocol: $0.944$ on
soft/general against LightGBM's $0.971$, and $-0.528$ on hard/safety
against LightGBM's $-0.024$. The additive structure with only ten
interaction terms is a genuine capacity restriction relative to a
fully-interacting tree ensemble.

\subsection{Why EBM is absent from the leave-rule-out results}

\begin{limitbox}[title=Stated gap in the revised evaluation]
EBM was \textbf{not} re-evaluated under leave-one-rule-out. The reason is
mechanical rather than scientific: the production EBM pipeline uses its own
preprocessing path (a frame-level imputer preserving native categorical
handling), not the shared scaled-and-one-hot \texttt{ColumnTransformer}
that every other model in the diagnostic harness uses. Running it inside
the shared harness would have changed its preprocessing as well as its
partitioning, confounding the comparison.

Consequently, \textbf{EBM's context-holdout numbers must not be placed
alongside the other models' leave-rule-out numbers as though they were
comparable.} They measure different things. Completing this evaluation
with EBM's own pipeline and the same rule-wise folds is a concrete,
bounded piece of outstanding work, listed in \Cref{ch:proposals}.
\end{limitbox}

\section{The alternative model families}

Thirteen further configurations were evaluated under both protocols. Their
behaviour under the protocol change is the clearest single illustration of
what the change measures.

\subsection{Regularised linear models}

Ridge, lasso, elastic net, Bayesian ridge and PLS regression all fit
\begin{equation}
\hat{y} = \beta_0 + \sum_{j} \beta_j x_j,
\end{equation}
differing only in the penalty applied to $\boldsymbol{\beta}$: $L_2$ for
ridge, $L_1$ for lasso, a convex combination for elastic net, a Gaussian
prior for Bayesian ridge, and a latent-component projection for PLS.

\begin{center}\small
\begin{tabular}{lrrrr}
\toprule
Model & \multicolumn{2}{c}{context-holdout $\rsq$} & \multicolumn{2}{c}{leave-rule-out $\rsq$} \\
\cmidrule(lr){2-3}\cmidrule(lr){4-5}
 & soft/general & mean, 6 spec. & mean & median \\
\midrule
ridge regression & 0.910 & 0.594 & $-0.874$ & $-0.798$ \\
lasso & 0.908 & 0.595 & $-0.690$ & $-0.736$ \\
elastic net & 0.895 & 0.592 & $-0.772$ & $-0.675$ \\
Bayesian ridge & 0.910 & 0.597 & $-0.872$ & $-0.795$ \\
PLS regression & 0.902 & 0.594 & $-1.106$ & $-0.952$ \\
\bottomrule
\end{tabular}
\end{center}

\noindent The two context-holdout columns are given separately because the
six-specialist mean is dragged down by the safety specialists, where every
model scores poorly; the soft/general column is the like-for-like
comparison against the production figure of $0.971$.

\begin{findingbox}[title=The linear-model collapse]
Ridge regression scores $\rsq = 0.910$ on soft/general under the original
protocol --- close enough to the production LightGBM's $0.971$ that it was
initially taken as evidence that model complexity was irrelevant to the
problem. Under leave-rule-out the same model scores a mean $\rsq$ of
$-0.874$: substantially worse than predicting a constant.

The interpretation is that a global linear fit can interpolate
\emph{within} the known rules, because within any one rule the mapping is
locally smooth --- but it has no mechanism for representing a rule it has
not seen, and its extrapolations are actively harmful. The conclusion from
Experiment 1 therefore needs a correction: model complexity is irrelevant
to \emph{within-rule interpolation} but decisive for \emph{cross-rule
transfer}.
\end{findingbox}

\subsection{Instance-based and kernel methods}

$k$-nearest neighbours ($k=5$) predicts the mean target of the five
nearest training rows in standardised feature space. It scores $0.905$ on
soft/general under the original protocol and a mean of $-3.412$ under
leave-rule-out --- the worst non-trivial model in the study. The failure
mode is transparent: when the entire mechanism is removed, the nearest
neighbours are drawn from different mechanisms, and their average carries
no relevant information.

RBF-kernel SVR fits
$\hat{y} = \sum_i \alpha_i \exp(-\gamma \lVert \mathbf{x} - \mathbf{x}_i \rVert^2) + b$,
with training subsampled to $4\,000$ rows for tractability. It scores
$0.946$ on soft/general then a mean of $-0.454$ under leave-rule-out: the
same pattern, less severe, because the kernel provides smoother
interpolation than a hard $k$-neighbour average.

\subsection{Tree ensembles outside the production set}

\begin{center}\small
\begin{tabular}{lrrr}
\toprule
Model & \multicolumn{2}{c}{context-holdout $\rsq$} & leave-rule-out \\
\cmidrule(lr){2-3}
 & soft/general & mean, 6 spec. & mean $\rsq$ \\
\midrule
extra trees & 0.961 & 0.608 & 0.209 \\
gradient boosting (sklearn) & 0.959 & 0.597 & 0.280 \\
AdaBoost & 0.917 & 0.601 & $-1.248$ \\
shallow random forest & 0.934 & 0.597 & $-0.230$ \\
shallow decision tree & 0.926 & 0.553 & $-1.733$ \\
small MLP & 0.972 & 0.595 & 0.219 \\
\bottomrule
\end{tabular}
\end{center}

Extra trees and sklearn gradient boosting remain positive under the
revised protocol --- they are structurally similar to the production
ensembles --- but at roughly one third of their performance. The
deliberately shallow variants collapse, confirming that depth is what
carries cross-rule transfer.

The small MLP is worth singling out: it \emph{matched} the production
LightGBM on soft/general under the original protocol ($0.972$ vs $0.971$)
and falls to a mean $\rsq$ of $0.219$ under leave-rule-out. It is the
sharpest single example of two models that an inadequate evaluation
declares equivalent and an adequate one separates by six-tenths of an
$\rsq$ point.

\subsection{The dummy baseline}

\texttt{DummyRegressor(strategy="mean")} predicts the training mean for
every row. It scores $-0.004$ under the original protocol --- correct,
since predicting the training mean on a test set drawn from the same
distribution yields approximately zero --- and $-15.637$ under
leave-rule-out, because each held-out rule's mean differs substantially
from the pooled training mean.

That number is a useful calibration for the panel: it establishes the
scale on which the leave-rule-out $\rsq$ values should be read. On that
scale, LightGBM's $+0.627$ and ridge's $-0.874$ are both far above the
no-information floor, and the separation between them is real.

\tbl{tab_exp3_extra_regressors}{All alternative models under the original
context-holdout protocol, averaged across the six specialists.}
